\documentclass[11pt,a4paper]{article}
\usepackage[margin=23mm]{geometry}
\usepackage{amsmath,amssymb,amsthm}
\setlength{\emergencystretch}{2em}
\newtheorem{lemma}{Lemma}
\newtheorem{theorem}[lemma]{Theorem}
\newcommand{\gen}[1]{\langle#1\rangle}
\title{Two blocks are incompatible with the critical head}
\author{Human Compression project: L044 continuation}
\date{}
\begin{document}\maketitle
Retain the complete critical-head hypotheses: $n=2r\ge18$,
$G\le\operatorname{Alt}(\Omega)$ is transitive, its string C-marking
has connected noncommutation diagram, and every full maximal omission
is intransitive. The actual first three marks form the acquired
six-point head on $D$. The acquired early-cut calculation gives
\[
 \tau=\rho_3=(6\ b)\delta,\quad \sigma=\rho_4,\quad
 x=b^\sigma\ne b,\quad \delta b=b,\quad\delta x=x,
\]
where $\delta$ is odd and supported on $E=\Omega\setminus D$.
The group $K=G_{\ge5}$ fixes $D\cup\{b\}$ and centralizes
$\delta$; $B=G_{\ge4}$ fixes $D$ pointwise.

\begin{lemma}[A saturated forest]
\label{forest}
Let $T=\gen{t_1,\ldots,t_{N-2}}$ act faithfully on $N$ points.
Assume the $t_i$ are nonidentity involutions, the noncommutation
 diagram on the distinguished marks is connected, $T$ is
intransitive, and deleting each $t_i$ changes its orbit partition.
Then $T$ fixes one point and is the full natural $S_{N-1}$ on
the others. Every distinguished mark is one transposition there.
\end{lemma}
\begin{proof}
For each $i$ select an edge $a_i\mathrel{-_{i}}a_i^{t_i}$
whose endpoints lie in different orbits of the omission of $t_i$.
These $N-2$ edges form a forest: a cycle would give an avoiding
path for one of its selected edges. Thus the selected forest has
exactly two components. Both lie in $T$-orbits. Since $T$ is
intransitive, these components are precisely its two orbits.

Deleting a selected edge $j$ splits its tree into two pieces.
The actual omission of $t_j$ contains all the other selected
edges, and cannot join these two pieces by eligibility. Hence
it preserves each piece, as well as the other tree component.
Every $t_i$ with $i\ne j$ preserves this cut. In a tree the
only two distinct vertices not separated by any cut other than
cut $i$ are the endpoints of edge $i$. It follows that $t_i$
is exactly its selected transposition. The transpositions of
each tree generate the full symmetric group on that tree.
If both trees had edges, their two nonempty collections of
marks would commute with one another, disconnecting the diagram.
Thus one tree is a singleton, as claimed.
\end{proof}

\begin{theorem}
Under the standing hypotheses, $B$ cannot be transitive on $E$
and preserve a partition of $E$ into two blocks of equal size.
\end{theorem}
\begin{proof}
Suppose that the blocks are $U,V$, each of size
$m=(n-6)/2=r-3\ge6$, with $b\in U$. Every mark of $K$
fixes $b$, so it preserves both blocks. Transitivity forces
$\sigma$ to interchange them. Being an involution, $\sigma$
is a matching of all $m$ points of $U$ with all $m$ points
of $V$. It is even and fixes $D$, so $m$ is even.
In particular $x\in V$.

Put $e=\rho_5$ and $H=G_{\ge6}$. Every element of $H$
preserves $U,V$, centralizes $\sigma$, and fixes $b,x$.
Identifying matched pairs by $u\leftrightarrow u^\sigma$,
$H$ acts diagonally on
\[
 U_0=U\setminus\{b\},\qquad V_0=V\setminus\{x\}.
\]
Its action $T$ on $U_0$ is faithful: a kernel element fixes
$U$, then fixes $V$ by commutation with $\sigma$, and fixes
$D$ already. Consequently the inherited string C-marking has
rank $r-6=m-3=N-2$ on $N=m-1$ points. Every mark is
nonidentity, and its noncommutation diagram is connected,
being a final interval of the original connected string.

Each row omission changes the row orbit partition. Otherwise
its identical diagonal action on $V_0$, together with the fixed
points $b,x,D$, would give the same point-orbit partition as
$H$. Adding the other distinguished generators would then give
the same orbit partition with or without the omitted generator:
orbits of a generated group are the equivalence relation generated
by the individual orbit relations. The full ambient omission
would be transitive, contrary to hypothesis.

The mark $e$ moves $x$. Indeed, if it fixed $x$, then
$\gen{e,H}$ would fix both $b,x$, while $\sigma$ interchanges
them, leaving $\{b,x\}$ invariant under $B$. This contradicts
transitivity since $|E|>2$. Write $y=x^e\in V_0$.
Since $\delta$ centralizes $e$ and fixes $x$, it fixes $y$.
It also centralizes $H$ and fixes $b,x$.

If $T$ is transitive, the accepted complete transitive-rank
theorem (L047, Theorem 1.2) applies to this faithful string
C-marking. Its degree $N=m-1\ge5$ is odd and its rank
$N-2>N/2$. Every proper exception in that theorem has even
degree. Thus $T$ is the natural $A_N$ or $S_N$.
The centralizing permutation $\delta$ fixes $y$, so it fixes
its entire $H$-orbit $V_0$ pointwise. It therefore preserves
$U_0$, the only other nontrivial $H$-orbit. The centralizer
of natural $A_N$ or $S_N$ in $S_N$ is trivial for $N\ge5$,
so $\delta$ fixes $U_0$ as well. Together with $b,x,D$,
these account for all points, contradicting oddness of $\delta$.

If $T$ is intransitive, Lemma~\ref{forest} applies. Let $p$
be its singleton in $U_0$ and $p'=p^\sigma$ the corresponding
singleton in $V_0$. On the remaining two rows, $H$ is the
diagonal natural $S_{m-2}$, where $m-2\ge4$. Its fixed
points in $E$ are precisely $b,x,p,p'$. On the two large
rows its centralizer consists of the identity and the full
matching that exchanges the rows: each separate natural
centralizer is trivial, and the matching is the unique
equivariant bijection. This exchange has $m-2$ transpositions,
an even number.

If $y=p'$, fixation of $b,x,y$ forces $\delta$ to fix all
four fixed points. Its remaining action is even, a contradiction.
Otherwise $y$ lies in a large row. Fixing $y$ excludes the
row exchange and makes $\delta$ fix both large rows pointwise.
As $\delta$ is odd and fixes $b,x$, it must be exactly
$(p\ p')$. But $H$ fixes this pair pointwise, $\sigma$
interchanges its two points, and $e$ preserves it because
$e$ centralizes $\delta$. Thus $B=\gen{\sigma,e,H}$
preserves $\{p,p'\}$, again contradicting transitivity.
Both cases are impossible.
\end{proof}

\paragraph{Attribution and boundary.}
The critical-head and early-cut calculations remain acquired
inputs. The complete L047 recovery, including its common theory,
finite boundary, software and distinct reviews, is charged as a
prior result of Cameron, Fernandes, Leemans and Mixer recovered
in this project. Only its odd-degree consequence is applied;
no L044 alternating rank induction is used here. The new argument
uses the diagonal row action and the odd exterior involution to
exclude two large blocks. Other block systems, opposite heads,
the unshared endpoint and the full L044 recovery remain separate.
No new finite census, formal theorem, whole-original or open-question
credit is claimed.
\end{document}
